The NRC Protein Folding Engine is founded on the principles of high-dimensional geometry and number theory. The engine utilizes the TUPT_E256 sequence over Z_12289, which provides a robust framework for representing biological sequences. The Mod 9 exclusion filters applied to the sequence ensure that invalid configurations are locked out, while the GTT Entropy target of ~10.96 nats guides the sequence towards biologically plausible conformations. The digital root audits, based on the TTT-7 theorem, provide an additional layer of validation for the predicted conformations. The digital root $dr$ is defined as the recursive sum of the digits of a number, modulo 9, and is used to identify sequences with desirable properties. The TTT-7 theorem states that a sequence with a digital root $dr \in \{1, 2, 4, 5, 7, 8\}$ is likely to exhibit biologically plausible behavior.

The mathematical foundations of the NRC Protein Folding Engine can be summarized as follows:

* High-dimensional manifold projections: The engine utilizes 2048D/8192D projections to represent biological sequences in a compact and efficient manner.
* Digital root audits: The engine applies the TTT-7 theorem to identify sequences with desirable properties, using the digital root $dr$ to guide the prediction process.
* Mod 9 exclusion filters: The engine applies Mod 9 exclusion filters to lock out invalid configurations and ensure that the predicted conformations are biologically plausible.